For fixed $B,C$, the locus of $D$ giving a right triangle is the circle with diameter $BC$ (right angle at $D$), together with the perpendicular lines to $BC$ through $B$ and $C$. Write $BC=a,CA=b,AB=c$, and $M$ for the midpoint of $BC$. The distance from $A$ to the circle is $A M - \dfrac{a}{2}$, since $A M > \dfrac{a}{2}$ in each of these acute side assignments. The median formula gives circle distances $\dfrac{\sqrt{145}-5}{2}$, $\dfrac{\sqrt{112}-6}{2}$, and $\dfrac{\sqrt{73}-7}{2}$ for $a=5,6,7$. The corresponding minimum distances to the two perpendicular lines are $\dfrac{6}{5} , 1 , \dfrac{18}{7}$, found from the projection $x = \dfrac{c^2+a^2-b^2}{2a}$ and $\min(x,a-x)$. The smallest of all six is $\dfrac{\sqrt{73}-7}{2}$. It is attained at the closest point of the diameter circle, which is not an endpoint. Therefore $\boxed{m+n+p=82}$.
\begin{asy}
import graph;
usepackage("amsmath");
import geometry;
size(7cm);
pair B=(-3.5,0),C=(3.5,0),A=(11/14,12sqrt(6)/7),M=(0,0),D=3.5A/abs(A);draw(B--A--C--cycle);draw(circle(M,3.5));draw(A--D,dashed);draw(B--D--C);label("$A$",A,(0,1));label("$B$",B,(-1,0));label("$C$",C,(1,0));label("$D$",D,(1,0));label("$M$",M,(0,-1));
\end{asy}

